\section{Discussion and Conclusion}

- Net effect of immigration vs population aging vs behavoir
- future expections about component and AI contributions
- Limites and room for improvement

Taken together, the two largest negative effects come from Spanish age structure, with young Spanish aged 16--34 and Spanish adults aged 35--54 contributing a combined \(-156.4\%\) of the total change. The third-largest negative effect comes from the age-structure component of young foreigners aged 16--34, at \(-34.2\%\). Against these downward pressures, the largest positive contribution comes from foreign population growth among adults aged 35--54, at \(86.3\%\). This is followed by higher activity rates among Spanish adults and older Spanish workers, which together contribute \(138.0\%\).

    Immigration is especially important through the population-size component: for foreigners, \textit{tpop} contributes about \(163.1\%\) of the total change. This implies that immigration, through foreign population growth, more than offset the negative age-structure effect associated with the ageing of the Spanish population, which amounts to about \(-116.8\%\).
    
    In short, the figure shows a clear contrast. Aggregate labour supply was supported by population growth, largely through immigration, and by higher participation among adults and older workers. At the same time, Spanish age-structure change and declining hours worked acted as major constraints. 



result summary
current labour force support policy
long term- encourage delaying retriement


- healthy aging improve labour market attachment of older workers 


Using data for people aged 50+ in 41 countries from 2000--2022, \textcite{gruss2025healthyAging} find that healthier ageing substantially improves older workers' labour-market outcomes. A decade of cognitive health gains raises labour-force participation by about 20 percentage points, weekly hours worked by around 6 hours, labour productivity by roughly 30\%, and annual labour earnings by about 35\%. Better health also delays retirement, increases weeks worked per year, and lowers unemployment risk. The main numeric implication is that healthier older cohorts can partly offset ageing-related declines in labour supply and productivity.


- Why did the foreign wpx and wph decreased between 2005 and 2025? 


Between 2005 and 2025, the foreign population increased substantially, but its average labour-market intensity declined. The foreign activity rate fell from \(74.77\%\) to \(71.44\%\), a decline of \(3.33\) percentage points, while average weekly hours decreased from \(42.69\) to \(40.11\), a decline of \(2.58\) hours per week. This raises an important question: why did participation and hours worked decline among foreigners while their population contribution increased?

A first explanation is compositional change within the foreign population. In the early 2000s, immigration to Spain was strongly connected to labour-market expansion and the arrival of working-age migrants \parencite{amuedo2008complements,arango2013exceptional,gonzalez2011very}. By 2025, the foreign population had become larger, more settled, and older. The share of foreigners aged 55+ increased from \(11.01\%\) to \(23.85\%\), a rise of \(12.84\) percentage points. This ageing of the foreign population mechanically lowers aggregate participation and average hours worked, consistent with the broader evidence that age composition affects aggregate labour supply \parencite{bde2023ageingParticipation,caixabank2017labourForceAgeing}.

A second explanation is exposure to labour-market volatility. Foreign workers were disproportionately affected by the 2008 financial crisis: immigrant employment declined, immigration inflows fell by more than 50\% between 2008 and 2013, and Spain experienced net migration outflows \parencite{rodriguez2016labor}. This suggests that foreign labour supply is more sensitive to economic cycles and employment instability. The decline in foreign \(wpx\) and \(whx\) therefore does not contradict the positive demographic contribution of immigration; rather, it shows that immigration expands potential labour supply mainly through population size and age structure, while participation and hours worked depend on labour-market integration, sectoral exposure, and employment stability \parencite{preston2014effect,fiorio2023migration,naumann2021population}.
